\documentclass[12pt]{article}

% Ћирилица: кодирање фонта и језик (српски, ћирилица)
\usepackage[T2A]{fontenc}
\usepackage[utf8]{inputenc}
\usepackage[serbianc]{babel}

% Величина стране и маргине
\usepackage[a4paper,top=2cm,bottom=2cm,left=2.5cm,right=2.5cm]{geometry}

% Корисни пакети
\usepackage{amsmath}
\usepackage{enumitem}
\usepackage{tcolorbox}
\tcbuselibrary{breakable}
\usepackage[colorlinks=true, allcolors=blue]{hyperref}

% Цртање електричних кола (европски симболи)
\usepackage[european]{circuitikz}
\usetikzlibrary{arrows.meta, decorations.markings, babel}

% Стрелица на средини линије (за струје у чвору)
\tikzset{
    sredstrelica/.style={postaction={decorate}, decoration={markings,
        mark=at position 0.55 with {\arrow{Stealth[length=7pt]}}}},
    obilazak/.style={-{Stealth[length=6pt]}, gray!70!black},
}

% Оквир за решене примере
\newtcolorbox[auto counter]{primer}{breakable,
    colback=white, colframe=black!55, colbacktitle=black!55, coltitle=white,
    fonttitle=\bfseries\small, title={Пример~\thetcbcounter},
    boxrule=0.5pt, arc=1.5pt, left=6pt, right=6pt, top=3pt, bottom=3pt,
    before skip=8pt, after skip=10pt}

% Оквир за кључне формуле
\newtcolorbox{kljucno}{colback=blue!4, colframe=blue!45!black,
    boxrule=0.5pt, arc=1.5pt, left=8pt, right=8pt, top=4pt, bottom=4pt}

\title{Кирхофови закони}
\author{Основи електротехнике 1}
\date{}

\begin{document}
\maketitle

%==================================================================
\section{Основни појмови}
%==================================================================

Пре него што уведемо Кирхофове законе, потребно је да знамо три појма:

\begin{itemize}
    \item \textbf{Чвор} је тачка у колу у којој се спајају најмање три
          проводника.
    \item \textbf{Грана} је део кола између два суседна чвора. Кроз све
          елементе једне гране тече иста струја.
    \item \textbf{Контура} је било који затворени пут у колу, који почиње и
          завршава се у истој тачки.
\end{itemize}

\begin{center}
\begin{circuitikz}
    \draw (0,0) to[battery1, l=$E$, invert] (0,3)
          to[R, l=$R_1$] (4,3)
          to[R, l_=$R_2$, -*] (4,0)
          -- (0,0);
    \draw (4,3) -- (7,3) to[R, l_=$R_3$] (7,0) -- (4,0);
    \fill (4,3) circle (2.5pt) node[above]{$A$};
    \node[below] at (4,0) {$B$};
\end{circuitikz}
\end{center}

У колу на слици:
\begin{itemize}
    \item тачке $A$ и $B$ су \textbf{чворови} (у свакој се спајају три
          проводника),
    \item постоје три \textbf{гране}: грана са извором $E$ и отпорником $R_1$,
          грана са отпорником $R_2$ и грана са отпорником $R_3$,
    \item једна \textbf{контура} је, на пример, затворени пут
          $E \to R_1 \to R_2 \to E$.
\end{itemize}

На шемама \textbf{дужа црта} у симболу извора означава позитиван пол ($+$),
а краћа негативан пол ($-$).

%==================================================================
\section{Први Кирхофов закон}
%==================================================================

Наелектрисање се у чвору не може ни створити ни изгубити. Зато колико
струје уђе у чвор, толико мора и да изађе из њега:

\begin{kljucno}
\[
\sum I_{\text{улазне}} = \sum I_{\text{излазне}}
\]
\end{kljucno}

Ово је као раскрсница цеви у водоводу: колико воде утекне у раскрсницу,
толико мора и да истекне из ње, јер се вода нигде не губи.

\begin{center}
\begin{tikzpicture}[thick]
    \draw[sredstrelica] (-2.2,0)    -- (0,0) node[pos=0.45, above]{$I_1$};
    \draw[sredstrelica] (0,2)       -- (0,0) node[pos=0.45, right]{$I_2$};
    \draw[sredstrelica] (0,0)       -- (2.2,0) node[pos=0.55, above]{$I_3$};
    \draw[sredstrelica] (0,0)       -- (0,-2) node[pos=0.55, right]{$I_4$};
    \fill (0,0) circle (3pt);
\end{tikzpicture}
\end{center}

За чвор на слици струје $I_1$ и $I_2$ улазе, а $I_3$ и $I_4$ излазе, па је:
\[
I_1 + I_2 = I_3 + I_4
\]

Закон се може написати и у облику $\sum I = 0$, ако струје које улазе у чвор
пишемо са знаком $+$, а струје које излазе са знаком $-$:
$I_1 + I_2 - I_3 - I_4 = 0$.

\textbf{Смер струје.} Када смер неке струје није познат, ми га
\textbf{претпоставимо} и нацртамо стрелицу. Ако на крају добијемо
\textbf{негативан} резултат, то значи да струја заправо тече у супротном смеру
од претпостављеног, а њена јачина је једнака добијеном броју без минуса.

\begin{primer}
У чвор на слици изнад улазе струје $I_1 = 5\,\text{A}$ и $I_2 = 3\,\text{A}$,
а излази струја $I_3 = 6\,\text{A}$. Израчунати струју $I_4$.

\textbf{Решење:}
\[
I_1 + I_2 = I_3 + I_4
\quad\Longrightarrow\quad
I_4 = I_1 + I_2 - I_3 = 5\,\text{A} + 3\,\text{A} - 6\,\text{A} = 2\,\text{A}
\]
\end{primer}

\begin{primer}
У чвор улази струја $I_1 = 2\,\text{A}$, а излази струја $I_2 = 7\,\text{A}$.
За трећу струју $I_3$ претпостављено је да \textbf{излази} из чвора.
Израчунати $I_3$ и одредити њен стварни смер.

\begin{center}
\begin{tikzpicture}[thick]
    \draw[sredstrelica] (-2.2,0) -- (0,0) node[pos=0.45, above]{$I_1$};
    \draw[sredstrelica] (0,0)    -- (2.2,0) node[pos=0.55, above]{$I_2$};
    \draw[sredstrelica] (0,0)    -- (0,-1.8) node[pos=0.55, right]{$I_3$};
    \fill (0,0) circle (3pt);
\end{tikzpicture}
\end{center}

\textbf{Решење:}
\[
I_1 = I_2 + I_3
\quad\Longrightarrow\quad
I_3 = I_1 - I_2 = 2\,\text{A} - 7\,\text{A} = -5\,\text{A}
\]
Резултат је негативан, па струја $I_3$ заправо \textbf{улази} у чвор и има
јачину $5\,\text{A}$.

\textbf{Провера:} у чвор улазе $2\,\text{A} + 5\,\text{A} = 7\,\text{A}$, а
излази $7\,\text{A}$.
\end{primer}

%==================================================================
\section{Други Кирхофов закон}
%==================================================================

Када обиђемо затворену контуру и вратимо се у почетну тачку, враћамо се на
исти потенцијал. Зато је у свакој затвореној контури збир електромоторних
сила (напона) једнак збиру падова напона на отпорницима:

\begin{kljucno}
\[
\sum E = \sum R \cdot I
\]
\end{kljucno}

\begin{kljucno}
\[
\sum V = \sum R \cdot I
\]
\end{kljucno}

Закон примењујемо овим редом:
\begin{enumerate}
    \item Изаберемо \textbf{смер обиласка} контуре (најчешће у смеру казаљке
          на сату) и означимо га кружном стрелицом.
    \item Електромоторну силу $E$ пишемо са знаком $+$ ако кроз извор
          пролазимо од $-$ ка $+$ пола, а са знаком $-$ ако пролазимо од $+$ ка
          $-$ полу.
    \item Пад напона $R \cdot I$ пишемо са знаком $+$ ако се смер струје
          поклапа са смером обиласка, а са знаком $-$ ако је супротан.
\end{enumerate}

\begin{primer}
У колу на слици је $E = 24\,\text{V}$, $R_1 = 4\,\Omega$ и
$R_2 = 8\,\Omega$. Израчунати струју у колу и напоне на отпорницима.

\begin{center}
\begin{circuitikz}
    \draw (0,0) to[battery1, l=$E$, invert] (0,3)
          to[short, i>^=$I$] (1,3)
          to[R, l=$R_1$] (4,3)
          to[R, l=$R_2$] (4,0)
          -- (0,0);
    \draw[obilazak] (2,2.1) arc[start angle=90, end angle=-200, radius=0.6];
\end{circuitikz}
\end{center}

\textbf{Решење:}

Обилазимо контуру у смеру казаљке на сату. Кроз извор пролазимо од $-$ ка
$+$, па је $E$ позитивна. Струја тече у смеру обиласка, па су оба пада
напона позитивна:
\[
E = R_1 I + R_2 I = (R_1 + R_2)\, I
\quad\Longrightarrow\quad
I = \frac{E}{R_1 + R_2} = \frac{24\,\text{V}}{4\,\Omega + 8\,\Omega} = 2\,\text{A}
\]
Напони на отпорницима:
\[
U_1 = R_1 I = 4\,\Omega \cdot 2\,\text{A} = 8\,\text{V}, \qquad
U_2 = R_2 I = 8\,\Omega \cdot 2\,\text{A} = 16\,\text{V}
\]
\textbf{Провера:} $U_1 + U_2 = 8\,\text{V} + 16\,\text{V} = 24\,\text{V} = E$.
\end{primer}

\begin{primer}
У колу на слици су два извора, $E_1 = 20\,\text{V}$ и $E_2 = 8\,\text{V}$, и
отпорници $R_1 = 2\,\Omega$ и $R_2 = 4\,\Omega$. Израчунати струју у колу.

\begin{center}
\begin{circuitikz}
    \draw (0,0) to[battery1, l=$E_1$, invert] (0,3)
          to[short, i>^=$I$] (1,3)
          to[R, l=$R_1$] (4,3)
          to[battery1, l=$E_2$] (4,0)
          to[R, l=$R_2$] (0,0);
    \draw[obilazak] (2,2.1) arc[start angle=90, end angle=-200, radius=0.6];
\end{circuitikz}
\end{center}

\textbf{Решење:}

Обилазимо контуру у смеру казаљке на сату:
\begin{itemize}
    \item кроз извор $E_1$ пролазимо од $-$ ка $+$, па је $E_1$ позитивна,
    \item кроз извор $E_2$ пролазимо од $+$ ка $-$, па је $E_2$ негативна,
    \item струја тече у смеру обиласка, па су падови напона позитивни.
\end{itemize}
\[
E_1 - E_2 = R_1 I + R_2 I
\quad\Longrightarrow\quad
I = \frac{E_1 - E_2}{R_1 + R_2}
  = \frac{20\,\text{V} - 8\,\text{V}}{2\,\Omega + 4\,\Omega}
  = \frac{12\,\text{V}}{6\,\Omega} = 2\,\text{A}
\]
Извори делују један против другог, па струју стварају само са разликом
својих електромоторних сила.
\end{primer}

%==================================================================
\section{Примена оба закона}
%==================================================================

У сложенијим колима користимо оба закона заједно: први закон за чворове, а
други за контуре.

\begin{primer}
Отпорници $R_1 = 6\,\Omega$ и $R_2 = 3\,\Omega$ везани су паралелно на извор
$E = 12\,\text{V}$. Помоћу Кирхофових закона израчунати струје $I_1$, $I_2$ и
$I$.

\begin{center}
\begin{circuitikz}
    \draw (0,0) to[battery1, l=$E$, invert] (0,3)
          to[short, i>^=$I$, -*] (3,3)
          to[R, l_=$R_1$, i>^=$I_1$, -*] (3,0)
          -- (0,0);
    \draw (3,3) -- (6,3)
          to[R, l_=$R_2$, i>^=$I_2$] (6,0)
          -- (3,0);
\end{circuitikz}
\end{center}

\textbf{Решење:}

\textbf{Други закон, контура} $E \to R_1$: у контури су само извор и
отпорник $R_1$:
\[
E = R_1 I_1
\quad\Longrightarrow\quad
I_1 = \frac{E}{R_1} = \frac{12\,\text{V}}{6\,\Omega} = 2\,\text{A}
\]

\textbf{Други закон, контура} $E \to R_2$:
\[
E = R_2 I_2
\quad\Longrightarrow\quad
I_2 = \frac{E}{R_2} = \frac{12\,\text{V}}{3\,\Omega} = 4\,\text{A}
\]

\textbf{Први закон, горњи чвор:} струја $I$ улази, а $I_1$ и $I_2$ излазе:
\[
I = I_1 + I_2 = 2\,\text{A} + 4\,\text{A} = 6\,\text{A}
\]
\end{primer}

\end{document}